Debido a la digitalización que se está produciendo en la realización de todo tipo de trámites a través de Internet, actualmente la sociedad tiene una dependencia creciente de las redes y sistemas informáticos. En este contexto, la ciberseguridad se ha convertido en un campo esencial para todo el mundo, tanto organizaciones como particulares.

En el ámbito de la ciberseguridad, los sistemas de detección de intrusiones tienen un papel importante, ya que analizan la actividad de una red o sistema con el objetivo de identificar comportamientos fuera de lo normal o maliciosos. No obstante, la evolución de los ataques hace que resulte necesario estudiar técnicas capaces de adaptarse a estos nuevos patrones de tráfico. Es aquí donde el aprendizaje automático representa una alternativa de interés, ya que permite la creación de modelos capaces de aprender constantemente y realizar predicciones sobre patrones que no se conocían previamente.